% -----------------------------------------------------------
\section{Atonement, Atone}
% -----------------------------------------------------------
	\label{ATONEMENT}
	
% ----------------------
\subsection{See also}
% ----------------------

	\hyperref[INTERCESSION]{Intercession}, \hyperref[RECONCILIATION]{Reconciliation}, \hyperref[SACRIFICE]{Sacrifice},

% ----------------------
\subsection{1828 American Dictionary of the English Language} % *1828
% ----------------------
	\label{ATONEMENT-1828}

	% -----
	\subsubsection{Atonement}
	% -----

		\begin{compactenum}
			\item Agreement; concord; reconciliation, after enmity or controversy.
			\item Expiation; satisfaction or reparation made by giving an equivalent for an injury, or by doing or suffering that which is received in satisfaction for an offense or injury; with for.
			\item In theology, the expiation of sin made by the obedience and personal sufferings of Christ.
		\end{compactenum}
		
	% -----
	\subsubsection{Atone}
	% -----
	
		% --
		\paragraph{Intransitive}
		% --
		
			\begin{compactenum}
				\item To agree; to be in accordance; to accord. [This sense is obsolete.]
				\item To stand as an equivalent; to make reparation, amends or satisfaction for an offense or a crime, by which reconciliation is procured between the offended and offending parties.
				\item To atone for, to make compensation or amends.
			\end{compactenum}
			
		% --
		\paragraph{Transitive}
		% --
		
			\begin{compactenum}
				\item To expiate; to answer or make satisfaction for.
				\item To reduce to concord; to reconcile, as parties at variance; to appease. [Not now used.]
			\end{compactenum}

% % ----------------------
% \subsection{Oxford English Dictionary} % *OED
% % ----------------------
	% \label{ATONEMENT-OED}

	% \begin{compactitem}
		% \item[]
	% \end{compactitem}

% ----------------------
\subsection{Scriptural Usage} % *SCRIPTURES
% ----------------------
	\label{ATONEMENT-SCRIPTURES}

	% -----
	\subsubsection{Old Testament} % *OT
	% -----
		\label{ATONEMENT-OT}
	
		% --
		\paragraph{KJV}
		% --
			\label{ATONEMENT-OT-KJV}
	
			\begin{tabular}{ll}
				Total occurrences in the OT (atonement, atone) & 81
			\end{tabular}
			
			\begin{compactdesc}
				\item[]
			\end{compactdesc}
			
		% % --
		% \paragraph{Hebrew Bible}
		% % --
			% \label{ATONEMENT-HEBREW}
		
			% %
			% \subparagraph{\texorpdfstring{\he{sample word}}{}}
			% %
			
		% % --
		% \paragraph{Septuagint}
		% % --
		
			% %
			% \subparagraph{\texorpdfstring{\gr{sample word}}{}}
			% %
		
	% -----
	\subsubsection{New Testament} % *NT
	% -----
		\label{ATONEMENT-NT}
	
		% --
		\paragraph{KJV}
		% --
			\label{ATONEMENT-NT-KJV}
	
			\begin{tabular}{ll}
				Total occurrences in the NT (atonement, atone) & 1
			\end{tabular}
			
			\begin{compactdesc}
				\item[{\hyperref[ROMANS-5-11]{Romans 5:11}}] And not only so, but we also joy in God through our Lord Jesus Christ, by whom we have now received the atonement.
			\end{compactdesc}
			
		% --
		\paragraph{Greek New Testament}
		% --
			\label{ATONEMENT-GREEK}
		
			%
			\subparagraph{\texorpdfstring{\gr{καταλλαγή}}{}}
			%
				\label{KATALLAGE}
			
				\begin{compactdesc}
					\item[BDAG 462a] \textbf{}
						\begin{compactitem}
							\item reestablishment of an interrupted or broken relationship, \textit{reconciliation}
						\end{compactitem}
				\end{compactdesc}
				
				\begin{compactdesc}
					\item[NT uses] \textbf{}
						\begin{compactdesc}
							\item[Romans 5:11] \gr{οὐ μόνον δέ, ἀλλὰ καὶ καυχώμενοι ἐν τῷ θεῷ διὰ τοῦ κυρίου ἡμῶν Ἰησοῦ Χριστοῦ, δι᾽ οὗ νῦν τὴν καταλλαγὴν ἐλάβομεν.}
							\item[Romans 11:15] \gr{εἰ γὰρ ἡ ἀποβολὴ αὐτῶν καταλλαγὴ κόσμου, τίς ἡ πρόσλημψις εἰ μὴ ζωὴ ἐκ νεκρῶν;}
							\item[2 Corinthians 5:18] \gr{τὰ δὲ πάντα ἐκ τοῦ θεοῦ τοῦ καταλλάξαντος ἡμᾶς ἑαυτῷ διὰ Χριστοῦ καὶ δόντος ἡμῖν τὴν διακονίαν τῆς καταλλαγῆς,}
							\item[2 Corinthians 5:19] \gr{ὡς ὅτι θεὸς ἦν ἐν Χριστῷ κόσμον καταλλάσσων ἑαυτῷ, μὴ λογιζόμενος αὐτοῖς τὰ παραπτώματα αὐτῶν, καὶ θέμενος ἐν ἡμῖν τὸν λόγον τῆς καταλλαγῆς.}
						\end{compactdesc}
				\end{compactdesc}
		
	% -----
	\subsubsection{Book of Mormon} % *BOM
	% -----
		\label{ATONEMENT-BOM}
	
		\begin{tabular}{ll}
			Total occurrences in the BOM (atonement, atone) & 39
		\end{tabular}
		
		\begin{compactdesc}
			\item[2 Nephi 2:10] And because of the intercession for all, all men come unto God; wherefore, they stand in the presence of him, to be judged of him according to the truth and holiness which is in him. Wherefore, the ends of the law which the Holy One hath given, unto the inflicting of the punishment which is affixed, which punishment that is affixed is in opposition to that of the happiness which is affixed, to answer the ends of the atonement---
			\item[2 Nephi 9:7] Wherefore, it must needs be an infinite atonement---save it should be an infinite atonement this corruption could not put on incorruption. Wherefore, the first judgment which came upon man must needs have remained to an endless duration. And if so, this flesh must have laid down to rot and to crumble to its mother earth, to rise no more.
			\item[2 Nephi 9:25] Wherefore, he has given a law; and where there is no law given there is no punishment; and where there is no punishment there is no condemnation; and where there is no condemnation the mercies of the Holy One of Israel have claim upon them, because of the atonement; for they are delivered by the power of him.
			\item[2 Nephi 9:26] For the atonement satisfieth the demands of his justice upon all those who have not the law given to them, that they are delivered from that awful monster, death and hell, and the devil, and the lake of fire and brimstone, which is endless torment; and they are restored to that God who gave them breath, which is the Holy One of Israel.
			\item[2 Nephi 10:25] Wherefore, may God raise you from death by the power of the resurrection, and also from everlasting death by the power of the atonement, that ye may be received into the eternal kingdom of God, that ye may praise him through grace divine. Amen.
			\item[2 Nephi 25:16] And after they have been scattered, and the Lord God hath scourged them by other nations for the space of many generations, yea, even down from generation to generation until they shall be persuaded to believe in Christ, the Son of God, and the atonement, which is infinite for all mankind---and when that day shall come that they shall believe in Christ, and worship the Father in his name, with pure hearts and clean hands, and look not forward any more for another Messiah, then, at that time, the day will come that it must needs be expedient that they should believe these things.
			\item[Jacob 4:11--12] Wherefore, beloved brethren, be reconciled unto him through the atonement of Christ, his Only Begotten Son, and ye may obtain a resurrection, according to the power of the resurrection which is in Christ, and be presented as the first-fruits of Christ unto God, having faith, and obtained a good hope of glory in him before he manifesteth himself in the flesh. 12 And now, beloved, marvel not that I tell you these things; for why not speak of the atonement of Christ, and attain to a perfect knowledge of him, as to attain to the knowledge of a resurrection and the world to come?
			\item[Jacob 7:12] And this is not all---it has been made manifest unto me, for I have heard and seen; and it also has been made manifest unto me by the power of the Holy Ghost; wherefore, I know if there should be no atonement made all mankind must be lost.
			\item[Mosiah 3:11] For behold, and also his blood atoneth for the sins of those who have fallen by the transgression of Adam, who have died not knowing the will of God concerning them, or who have ignorantly sinned.
			\item[Mosiah 3:15] nd many signs, and wonders, and types, and shadows showed he unto them, concerning his coming; and also holy prophets spake unto them concerning his coming; and yet they hardened their hearts, and understood not that the law of Moses availeth nothing except it were through the atonement of his blood.
			\item[Mosiah 3:16] And even if it were possible that little children could sin they could not be saved; but I say unto you they are blessed; for behold, as in Adam, or by nature, they fall, even so the blood of Christ atoneth for their sins.
			\item[Mosiah 3:18] For behold he judgeth, and his judgment is just; and the infant perisheth not that dieth in his infancy; but men drink damnation to their own souls except they humble themselves and become as little children, and believe that salvation was, and is, and is to come, in and through the atoning blood of Christ, the Lord Omnipotent.
			\item[Mosiah 3:19] For the natural man is an enemy to God, and has been from the fall of Adam, and will be, forever and ever, unless he yields to the enticings of the Holy Spirit, and putteth off the natural man and becometh a saint through the atonement of Christ the Lord, and becometh as a child, submissive, meek, humble, patient, full of love, willing to submit to all things which the Lord seeth fit to inflict upon him, even as a child doth submit to his father.
			\item[Mosiah 4:2] And they had viewed themselves in their own carnal state, even less than the dust of the earth.  And they all cried aloud with one voice, saying: O have mercy, and apply the atoning blood of Christ that we may receive forgiveness of our sins, and our hearts may be purified; for we believe in Jesus Christ, the Son of God, who created heaven and earth, and all things; who shall come down among the children of men.
			\item[Mosiah 4:6] I say unto you, if ye have come to a knowledge of the goodness of God, and his matchless power, and his wisdom, and his patience, and his long-suffering towards the children of men; and also, the atonement which has been prepared from the foundation of the world, that thereby salvation might come to him that should put his trust in the Lord, and should be diligent in keeping his commandments, and continue in the faith even unto the end of his life, I mean the life of the mortal body---
			\item[Mosiah 4:7] I say, that this is the man who receiveth salvation, through the atonement which was prepared from the foundation of the world for all mankind, which ever were since the fall of Adam, or who are, or who ever shall be, even unto the end of the world.
		\end{compactdesc}
		
	% -----
	\subsubsection{Doctrine and Covenants} % *DC
	% -----
		\label{ATONEMENT-DC}
	
		\begin{tabular}{ll}
			Total occurrences in the DC & 
		\end{tabular}
		
		\begin{compactdesc}
			\item[]
		\end{compactdesc}
		
	% -----
	\subsubsection{Pearl of Great Price} % *PGP
	% -----
		\label{ATONEMENT-PGP}
	
		\begin{tabular}{ll}
			Total occurrences in the PGP & 
		\end{tabular}
		
		\begin{compactdesc}
			\item[]
		\end{compactdesc}
		
% ----------------------
\subsection{Key Phrases} % *KEYPHRASES *PHRASES
% ----------------------
	\label{ATONEMENT-PHRASES}

	\begin{compactdesc}
		\item[infinite atonement] \textbf{4} | \hyperref[2NEPHI-9-7]{2 Nephi 9:7} (2), \hyperref[2NEPHI-25-16]{25:16}; \hyperref[ALMA-34-12]{Alma 34:12}
			\begin{compactdesc}
				\item[Commentary] See also Alma 34:8--16, which speaks of an ``infinite and eternal'' sacrifice.
					\begin{compactitem}
						\item See the \hyperref[ATONEMENT-KEYS]{key points} below, about how \hyperref[2NEPHI-9-21]{2 Nephi 9:21} appears to put a restriction on it. ``Infinite'' probably can't mean ``unbound in every respect.'' I'm not even sure that would make sense. 
					\end{compactitem}
			\end{compactdesc}
		\item[perfect atonement] \textbf{1} | \hyperref[DC76-69]{DC 76:69}
		\item[power of (the) atonement] 
		\item[through the atonement] \textbf{8} | \hyperref[JACOB-4-11]{Jacob 4:11}; \hyperref[MOSIAH-3-15]{Mosiah 3:15}, \hyperref[MOSIAH-3-19]{19}, \hyperref[MOSIAH-4-7]{4:7}; \hyperref[ALMA-13-5]{Alma 13:5}, \hyperref[ALMA-34-9]{34:9}; \hyperref[MORONI-7-41]{Moroni 7:41}; \hyperref[DC74-7]{DC 74:7}
			\begin{compactdesc}
				\item[Commentary] There is technically a ninth occurrence in the \hyperref[ARTICLESOFFAITH-3]{third article of faith}. 
			\end{compactdesc}
		\item[through the atoning blood] \textbf{2} | \hyperref[MOSIAH-3-18]{Mosiah 3:18}; \hyperref[HELAMAN-5-9]{Helaman 5:9}
			\begin{compactdesc}
				\item[Commentary] See also \hyperref[MOSIAH-3-15]{Mosiah 3:15}, which has ``through the atonement of his blood.''
			\end{compactdesc}
	\end{compactdesc}
	
% % ----------------------
% \subsection{Bible Dictionaries} % *DICTIONARIES
% % ----------------------
	% \label{ATONEMENT-BD}

% % ----------------------
% \subsection{Restoration Usage} % *RESTORATION
% % ----------------------
	% \label{ATONEMENT-RESTORATION}

	% % -----
	% \subsubsection{Joseph Smith} % JS
	% % -----
		% \label{ATONEMENT-JS}
	
		% \begin{compactdesc}
			% \item[DATE/SOURCE]
		% \end{compactdesc}
	
	% % -----
	% \subsubsection{Brigham Young} % BY
	% % -----
		% \label{ATONEMENT-BY}
	
		% \begin{compactdesc}
			% \item[DATE/SOURCE]
		% \end{compactdesc}
	
% ----------------------
\subsection{Commentary and Studies} % *COMMENTARY *STUDIES
% ----------------------
	\label{ATONEMENT-COMMENTARY}
	
	% -----
	\subsubsection{Key Points}
	% -----
		\label{ATONEMENT-KEYS}
	
		\begin{compactitem}
			\item General list of scriptures that refer to different aspects of the atonement
				\begin{compactitem}
					\item \hyperref[ISAIAH-53-3]{Isaiah 53:3--5}
					\item \hyperref[ROMANS-3-25]{Romans 3:25--26}
						\begin{compactitem}
							\item Notice that the emphasis is on God's righteousness here, similar to the verses at the beginning of DC 45.
						\end{compactitem}
					\item \hyperref[GALATIANS-2-19]{Galatians 2:19--21}
					\item \hyperref[GALATIANS-3-13]{Galatians 3:13}
					\item \hyperref[EPHESIANS-1-7]{Ephesians 1:7}
					\item \hyperref[EPHESIANS-2-13]{Ephesians 2:13--22}
					\item \hyperref[2CORINTHIANS-5-14]{2 Corinthians 5:14--21}
					\item \hyperref[COLOSSIANS-1-9]{Colossians 1:9--29}
					\item \hyperref[1NEPHI-12-10]{1 Nephi 12:10--11}
					\item \hyperref[1NEPHI-19-8]{1 Nephi 19:8---12}
					\item \hyperref[2NEPHI-9]{2 Nephi 9} 
						\begin{compactitem}
							\item Many things here; you'd have to go through the chapter to list them individually.
						\end{compactitem}
					\item \hyperref[2NEPHI-10]{2 Nephi 10} 
						\begin{compactitem}
							\item The stuff at the end specifically
						\end{compactitem}
					\item \hyperref[2NEPHI-25-16]{2 Nephi 25:16}
						\begin{compactitem}
							\item Mentions that the atonement is infinite; good comparison to Alma 34.
						\end{compactitem}
					\item \hyperref[2NEPHI-25-23]{2 Nephi 25:23}
					\item \hyperref[JACOB-4-11]{Jacob 4:11--12}
					\item \hyperref[MOSIAH-2-38]{Mosiah 2:38--39}; important to read this one in the light of other verses that use the term ``demands of justice,'' because \textit{mercy} is what satisfies those demands
					\item \hyperref[MOSIAH-3-7]{Mosiah 3:7}, \hyperref[MOSIAH-3-11]{11}, \hyperref[MOSIAH-3-13]{13}, \hyperref[MOSIAH-3-16]{16--18}, \hyperref[MOSIAH-3-21]{21}, \hyperref[MOSIAH-3-25]{25--26}
					\item \hyperref[MOSIAH-4-2]{Mosiah 4:2--3}
					\item \hyperref[MOSIAH-4-6]{Mosiah 4:6--8}
					\item \hyperref[MOSIAH-5-2]{Mosiah 5:2--15}
					\item \hyperref[MOSIAH-13-27]{Mosiah 13:27--28}
					\item \hyperref[MOSIAH-14]{Mosiah 14--16}
						\begin{compactitem}
							\item There's a huge amount in both of these chapters on the atonement, in a bunch of different ways.
						\end{compactitem}
					\item \hyperref[MOSIAH-18-2]{Mosiah 18:2}
					\item \hyperref[MOSIAH-26-22]{Mosiah 26:22--24}
					\item \hyperref[ALMA-11-40]{Alma 11:40--41}
					\item \hyperref[ALMA-22-13]{Alma 22:13--14}
					\item \hyperref[ALMA-24-10]{Alma 24:10--15}
					\item \hyperref[ALMA-33-22]{Alma 33:22}
					\item \hyperref[ALMA-34-8]{Alma 34:8--18}
					\item \hyperref[HELAMAN-5-9]{Helaman 5:9--12}
					\item \hyperref[HELAMAN-8-14]{Helaman 8:14--15}
					\item \hyperref[HELAMAN-14]{Helaman 14}
						\begin{compactitem}
							\item Most of the chapter, but not all of it.
							\item Through to verse 19 at the start.
							\item 28--31 at the end maybe?
						\end{compactitem}
					\item \hyperref[3NEPHI-11]{3 Nephi 11}
					\item \hyperref[MORMON-9-12]{Mormon 9:12--14}
					\item \hyperref[MORONI-6-4]{Moroni 6:4}, specifically I'm thinking of the ``author and finisher'' part
					\item \hyperref[MORONI-7-27]{Moroni 7:27--28}
					\item \hyperref[MORONI-7-41]{Moroni 7:41}
					\item \hyperref[MORONI-8-19]{Moroni 8:19--23}
					\item \hyperref[MORONI-10-33]{Moroni 10:33}
					\item \hyperref[DC18-11]{DC 18:11--12}
					\item \hyperref[DC19-15]{DC 19:15--19}
					\item \hyperref[DC20-19]{DC 20:19--28}
					\item \hyperref[DC21-9]{DC 21:9}
					\item \hyperref[DC35-2]{DC 35:2}
					\item \hyperref[DC38-4]{DC 38:4}
					\item \hyperref[DC45-3]{DC 45:3--5}
					\item \hyperref[DC45-51]{DC 45:51--52}
					\item \hyperref[DC53-2]{DC 53:2}
					\item \hyperref[DC54-1]{DC 54:1}
					\item \hyperref[DC76-69]{DC 76:69}
				\end{compactitem}
			\item Christ suffers ``the pains of all men, yea, the pains of every living creature, both men women and children, which belong to the family of Adam'' (\hyperref[2NEPHI-9-21]{2 Nephi 9:21}).
				\begin{compactitem}
					\item Note that this \textit{restricts} the atonement, or at least the pains suffered, to those beings who belong to the family of Adam. So it isn't ``infinite'' in this sense; we have to qualify it in some way.
					\item See also \hyperref[DC18-11]{DC 18:11--12} here.
				\end{compactitem}
		\end{compactitem}
	
	% -----
	\subsubsection{Categories of people the atonement applies to}
	% -----
		\label{ATONEMENT-categories}
	
		The power of the atonement is effective in the lives of many people for many different reasons. Here are some of the categories of people to which it applies, and the way in which it applies:
		
			\begin{enumerate}
				\item \textbf{Those who have fallen by the transgression of Adam.} See \hyperref[MOSIAH-3-11]{Mosiah 3:11}. ``His blood atoneth for the sins of'' these people. See more about this idea of the ``transgression of Adam'' \hyperref[TRANSGRESSION-PHRASES]{here}.
				\item \textbf{Those who have died not knowing the will of God concerning them.} See \hyperref[MOSIAH-3-11]{Mosiah 3:11}. ``His blood atoneth for the sins of'' these people. See also \hyperref[3NEPHI-6-18]{3 Nephi 6:18} and \hyperref[NUMBERS-15-27]{Numbers 15:27--29}.
				\item \textbf{Those who have ignorantly sinned.} See \hyperref[MOSIAH-3-11]{Mosiah 3:11}. ``His blood atoneth for the sins of'' these people. See also \hyperref[ROMANS-7-9]{Romans 7:9} and \hyperref[DC76-72]{DC 76:72}.
				\item \textbf{Little children.} See \hyperref[MOSIAH-3-16]{Mosiah 3:16}. ``The blood of Christ atones for their sins.'' See also \hyperref[DC29-46]{DC 29:46--47}, \hyperref[DC74-7]{DC 74:7}, 
			\end{enumerate}
			
	% -----
	\subsubsection{Temporal nature of the atonement}
	% -----
	
		See the note at \hyperref[MOSIAH-3-13]{Mosiah 3:13}.
		
		\begin{compactitem}
			\item See also \hyperref[JACOB-4-4]{Jacob 4:4}, \hyperref[JAROM-1-11]{Jarom 1:11}, \hyperref[MOSIAH-16-6]{Mosiah 16:6},
		\end{compactitem}
		
	% -----
	\subsubsection{Will we have to perform an atonement?}
	% -----
	
		\begin{compactitem}
			\item \hyperref[ROMANS-12-1]{Romans 12:1}, with the idea of our bodies as a ``living sacrifice''
		\end{compactitem}